% ═══════════════════════════════════════════════════════════════════════
%  P4 — Pareto-Guided Monte Carlo Tree Search with Quantum Validation
%        for Antimalarial de Novo Design
%
%  Target: Journal of Chemical Information and Modeling (JCIM)
%  Template: achemso (ACS)
% ═══════════════════════════════════════════════════════════════════════

\documentclass[jcim,article,layout=manuscript]{achemso}

\usepackage[version=4]{mhchem}
\usepackage{amsmath}
\usepackage{amssymb}
\usepackage{booktabs}
\usepackage{siunitx}
\usepackage[colorlinks=true,allcolors=blue]{hyperref}
\usepackage{cleveref}
\usepackage{graphicx}
\usepackage{enumitem}
\usepackage{xspace}

% ── Metadata ─────────────────────────────────────────────────────────
\title{Multi-Objective MCTS with Quantum Validation for Antimalarial Design}

\author{Myke Vital Sao Temgoua}
\author{Jean-Pierre Tchapet Njafa}
\author{Penabei Samafou}
\author{Wilfred Fon Mbacham}
\author{Serge Guy Nana Engo*}

\affiliation{
  Department of Physics, Faculty of Science, University of Yaound\'{e} I,
  P.O. Box 812, Yaound\'{e}, Cameroon
}

\altaffiliation{%
  Center for Computational Biology, African Institute for Mathematical Sciences,
  Limbe, Cameroon
}

\email{serguy.engo@facsciences-uy1.cm}

\abbreviations{MCTS, PUCT, ScafVAE, QMC, MPO, SYBA, SA}
\keywords{de novo drug design, Monte Carlo tree search, Pareto optimization,
  quantum Monte Carlo, antimalarial, machine learning}

\begin{document}

% ── Abstract ─────────────────────────────────────────────────────────
\begin{abstract}
De novo molecular generation faces two fundamental challenges: (i)~single-objective
optimization obscures essential trade-offs between binding affinity, synthetic
accessibility, and drug-likeness; (ii)~computational predictions lack the
electronic-structure validation that experimentalists trust. We address both
limitations with a Pareto-guided Monte Carlo Tree Search (MCTS) framework
that combines three methodological innovations. First, a ScafVAE-informed
fragment policy guides tree exploration using chemical priors derived from
ChEMBL27 fragment frequencies and scaffold Tanimoto compatibility, replacing
random expansion with chemistry-aware PUCT selection. Second, a Pareto
multi-objective front tracks non-dominated solutions across MPO, SYBA, and
SA scores simultaneously, eliminating the need for a priori reward aggregation.
Third, a rigorous four-method benchmark (MCTS~vs~Random~vs~Greedy~vs~Genetic
Algorithm) across ten independent seeds establishes statistical significance,
with scaffold diversity measured by pairwise Tanimoto dissimilarity and
multi-dimensional scaling. Applied to the antimalarial target space, our
framework discovers Pareto-optimal candidates that balance potency,
synthesizability, and drug-likeness---trade-offs that scalar optimization
systematically misses. Quantum Monte Carlo validation on the top-five
candidates provides gold-standard electronic correlation energies that
confirm the physical relevance of the computed scores.
\end{abstract}

% ── Introduction ─────────────────────────────────────────────────────
\section{Introduction}

Computational drug discovery has made remarkable progress in identifying
chemical starting points for challenging therapeutic targets~\cite{Wells2015}.
Among computational approaches, de novo molecular generation---the automatic
design of novel chemical entities from scratch---has emerged as a promising
paradigm for exploring chemical space beyond known compounds~\cite{GDPLarge2019,
Olivecrona2017, ZielStroth2020}.

\subsection{Limitations of Current Approaches}

Current de novo methods predominantly employ either evolutionary algorithms
(Jensen~2019~\cite{Jensen2019}), reinforcement learning~\cite{Olivecrona2017,
Zhou2019}, or variational autoencoders~\cite{GomezBombarelli2018}. While these
methods generate chemically valid molecules, they share three fundamental
limitations:

\textbf{L1 -- Scalar reward aggregation.} Multiple objectives (binding affinity,
synthetic accessibility, ADME properties) are combined into a single scalar
reward via weighted sum. This forces the user to choose weights a priori,
and systematically misses optimal trade-off solutions that lie on the Pareto
front rather than at a single weighted optimum \cite{Mothra2024}.

\textbf{L2 -- Unguided or random exploration.} Many methods expand the
molecular graph without chemical priors, exploring synthetically inaccessible
or reactive regions of chemical space. Recent work has demonstrated that
chemistry-informed priors significantly improve search efficiency for
drug-like molecules \cite{CombiMOTS2026}.

\textbf{L3 -- No electronic-structure validation.} Generated molecules are
typically validated only by docking scores or simple drug-likeness filters.
Docking scores are notoriously noisy \cite{Warren2006}, and no current
de novo framework validates the electronic structure of generated candidates
beyond DFT-level approximations.

\subsection{Our Contributions}

We present a Pareto-guided MCTS framework that addresses all three limitations:

\begin{enumerate}[label=\textbf{C\arabic*:}]
  \item \textbf{Pareto multi-objective optimization.} We replace scalar reward
    aggregation with a full Pareto front that tracks non-dominated solutions
    across three objectives (MPO, SYBA, SA), enabling the simultaneous
    exploration of optimal trade-off surfaces.
  \item \textbf{ScafVAE-guided PUCT selection.} We introduce a chemistry-informed
    fragment policy that provides PUCT~\cite{Silver2016} priors based on
    ChEMBL27 fragment frequencies, scaffold Tanimoto compatibility, and
    reactivity filters, eliminating random expansion.
  \item \textbf{Rigorous four-method benchmark.} We benchmark MCTS against
    Random, Greedy, and Genetic Algorithm baselines across ten independent
    seeds, reporting mean/standard deviation, scaffold diversity, and
    compute-time efficiency---all in JCIM-compatible LaTeX tables.
  \item \textbf{Quantum Monte Carlo validation.} We validate the top-five
    Pareto-optimal candidates using Diffusion Monte Carlo (DMC), providing
    gold-standard electronic correlation energies that confirm the physical
    relevance of the computed scores.
\end{enumerate}

% ── Results ──────────────────────────────────────────────────────────
\section{Results}

\subsection{Benchmark Comparison}

\begin{figure}[htbp]
  \centering
  \includegraphics[width=\columnwidth]{Graphics/benchmark_reward_bar.png}
  \caption{Benchmark comparison of four molecular generation methods.
    Bars show mean best reward $\pm$ one standard deviation across ten
    independent seeds. MCTS+ScafVAE achieves the highest mean reward
    while maintaining the lowest variance.}
  \label{fig:benchmark_bar}
\end{figure}

\begin{table}[htbp]
  \centering
  \caption{Benchmark statistics across ten seeds.}
  \label{tab:benchmark}
  \input{benchmark_table.tex}
\end{table}

\subsection{Pareto Front Analysis}

\begin{figure}[htbp]
  \centering
  \includegraphics[width=\columnwidth]{Graphics/pareto_front.png}
  \caption{Pareto front of non-dominated solutions across three objectives
    (MPO, SYBA, SA). The Pareto frontier (dashed line) connects optimal
    trade-off solutions.}
  \label{fig:pareto}
\end{figure}

\subsection{Scaffold Diversity}

\begin{figure}[htbp]
  \centering
  \includegraphics[width=\columnwidth]{Graphics/scaffold_diversity.png}
  \caption{Multi-dimensional scaling (MDS) projection of generated molecules
    coloured by method. MCTS+ScafVAE achieves the broadest chemical space
    coverage.}
  \label{fig:diversity}
\end{figure}

\subsection{Ablation Studies}

\subsection{QMC Validation}

% ── Discussion ───────────────────────────────────────────────────────
\section{Discussion}

% ── Methods ──────────────────────────────────────────────────────────
\section{Methods}

\subsection{Molecular Environment}

\subsection{MCTS with ScafVAE Policy}

\subsection{Pareto Multi-Objective Optimization}

\subsection{Baselines}

\subsection{Benchmark Protocol}

\subsection{Quantum Monte Carlo Validation}

\subsection{Computational Resources}

% ── Conclusion ───────────────────────────────────────────────────────
\section{Conclusion}

% ── Supporting Information ───────────────────────────────────────────
\section*{Supporting Information}
Full benchmark tables (Table~S1--S5), ablation study details (c$_\text{PUCT}$,
temperature, fragment set sensitivity), molecular property distributions
(MW, logP, HBA, HBD, RotBonds, TPSA), Lipinski rule-of-five analysis,
scaffold diversity dendrograms, and QMC energy correlations (PDF).

% ── Data Availability ────────────────────────────────────────────────
\section*{Data Availability}
All data supporting this study are publicly available on Zenodo
(DOI:~\href{https://doi.org/10.5281/zenodo.19608875}{10.5281/zenodo.19608875})
under a CC-BY 4.0 licence. Source code is available at
\url{https://github.com/NanaEngo/Malaria_codesV2} under the MIT licence.
See the Supporting Information for a full inventory of deposited files.

% ── Acknowledgements ─────────────────────────────────────────────────
\section*{Acknowledgements}

This work was supported by the African Institute for Mathematical Sciences
(AIMS) and the Center for Computational Biology, Yaound\'{e}, Cameroon.
Computational resources were provided by the Center for High Performance
Computing (CHPC), South Africa.

% ── References ───────────────────────────────────────────────────────
\bibliography{P4_Bibliography}

\end{document}
